\begin{abstract}

Deploying reinforcement learning agents on resource-constrained hardware requires compressing large teacher policies into compact students via behavioral cloning.
A practitioner must choose a student architecture---GRU, LSTM, Transformer, or spectral state space models---but it is unclear how much this choice matters relative to the quality of the teacher's demonstrations.

We study this question systematically, comparing four temporal architectures and two feedforward baselines across synthetic memory tasks, MuJoCo locomotion POMDPs, and MetaDrive autonomous driving, with over 500 experiments.
Our central finding is that a better teacher helps far more than a fancier architecture.
In a variance decomposition across all experiments, how good the teacher is explains $49$--$72\%$ of the difference in student performance, while architecture choice explains only $5$--$14\%$.
But architecture is not irrelevant: when the teacher is bad, architecture matters much more ($17$--$25\%$ interaction effect)---so the one time you most need to pick the right architecture is exactly when your training data is worst.

Architecture does matter in specific ways.
Recurrent models solve retention tasks perfectly; Transformers dominate content-addressable recall; and the best architecture for temporal inference depends on the task---GRU leads on MuJoCo locomotion, but STU and Transformer lead on MetaDrive autonomous driving, where lidar-rich observations favor spatial processing over compact recurrence.
A parameter-matched comparison confirms GRU's locomotion advantage is architectural, not capacity-driven.
We also document that the Flash STU tensordot approximation fails catastrophically at compact model scales, a finding unreported in prior work.

All code and results: \url{https://github.com/wasumayan/wasumind}.

\end{abstract}
